% RH_MODULE_v1_4.tex: route H for strips at 20 <= J <= 300, a module for Section 6 of the next paper version. It replaces the paragraph
% "The cover of [20,300]" of G2_PAPER_DRAFT_v4_12 (l.2067-2074) and stands before "Strips at strong coupling". Labels carry the prefix rh.
% v1_1: the fixes of the audit (rv01/AUDIT_RH_MODULE_v1_rv01.md, 1a77d790): M1, M2, c1-c6 and the S1 status comments. v1 (209ee05b)
% stays frozen.
% v1_2: SOURCE numbers from the final ru2 run over [20, 300] (route H CLOSED; the approval logged 2026-11-07T11:26:55Z) and the optional
% kappa >= 1 clause in the proof of Lemma rhslack (rv01/AUDIT_RH_MODULE_v1_rv01_ADD1.md fc400fc4, M1 row). v1_1 (1b1470fc) stays frozen.
% v1_3: comment lines only. The record of the final ru2 run is the frozen output pk05/out/ru2_U2.out with its manifest
% pk05/MANIFEST_ru2_U2_out.sha256; the log pk05/RU2_LOG_pk05.md is a live file, cited by its lines l.48-52 (
% AUDIT_G2KERNEL_LA_v1_1_DELTA_rv05_ADD10.md 715ff955, item m1; the note logged 2026-11-07T11:52:05Z). v1_2 (245b1dcb) stays frozen.
% v1_4 (pk14): route H's cover at N = 4 extends to [20, 350] (pk08/R11_ADDENDUM_N4_300_350_pk08_v1_1.md 3797ca7e
% CHECK_R11_ADDENDUM_N4_300_350_rv01.md d86b2e0f): the N = 4 range in the text, in Theorem thm:rhcover and in Corollary cor:rhstrips, and the
% SOURCE of the further ru2 run. The S8 batch logged 2026-11-08T13:58:59Z in pk14/NOTES_pk14.md. v1_3 (859366b0) stays frozen.
% SOURCE (notes, frozen): pk05/LEMMA_STRIP_RH_v1_3.md (3dab1187), sections 0, R1-R6, F and 7; pk05/LEMMA_STRIP_R7R8_v1_1.md (aa544045);
% pk05/LEMMA_STRIP_R9_v1_4.md (9f7f5905); pk05/LEMMA_STRIP_R10_v1_4.md (bd7f8861); pk05/LEMMA_STRIP_R11_v1_4.md (5e532179).
% SOURCE (reviews, no open item): rv01/REVIEW_LEMMA_STRIP_RH_v1_2_rv01_ADD1.md (1b78a0cf), REVIEW_LEMMA_STRIP_R7R8_v1_rv01_ADD1.md
% (2935195f), REVIEW_LEMMA_STRIP_R9_v1_rv01_ADD3.md (b3a11d86), REVIEW_LEMMA_STRIP_R10_v1_1_rv01_ADD1.md (ba194e9c),
% REVIEW_LEMMA_STRIP_R11_v1_1_rv01_ADD1.md (f6783f80); rv03/REVIEW_R2_STRIP_rv03_v1.md (eee9c197), R2 CONFIRMED.
% SOURCE (code): engine pk05/cs33h_1_rh_cert.py (b0f40c89; MANIFEST_cs33h_1_rh.sha256 df2233ff) with its core cs33b_rh_core.py (929ab8a0)
% and hubcert_v1e.py (c08d33c1), both loaded read-only after sha checks; driver pd5_driver.sh (9cb4e05d); union check ru2_ledger_union.py
% (837a50b3; MANIFEST_ru2.sha256 5d1c19ab).
% Notation against the notes: a_l, g_l, beta are k_l, r_l, J here; the vertical insertion X is U_e; J(p), I(e,p), I'(e,p) are w_p, w_{e,p},
% w'_{e,p}; beta_a is theta_a; y_e is z_e; slack_f is eps_sl; pi(t) is varpi(t); the modes are chi_0..chi_5 with nu_0 > ... > nu_5 as in R7R8.
% Wording of the coverage statement: the decision of 2026-11-06T22:48:58Z (pk13/NOTES_pk13.md l.1856): Cov >= 0 for every pair at
% J > 20, strict for the pairs the certificates cover; pairs with a common vertex take >= 0 from Theorem thm:shared.

\subsection{Strips for $20\le J\le300$}\label{sec:rh}

In this subsection $N\in\{3,4\}$ and all couplings of $P_3\times P_n$ equal $J$. We write $\sigma_1,\sigma_2,\sigma_3$ for the top, middle and bottom
spins of a column. The horizontal edges of the three rows have the types th, mh and bh, and the vertical edges the types tv (rows $1,2$) and bv
(rows $2,3$). We prove $\Cov(s_e,s_f)>0$ for $20\le J\le300$ ($20\le J\le350$ when $N=4$), every $n\ge4$ and every pair of distinct edges other than the two vertical edges of one
column (Theorem~\ref{thm:rhcover}). As in Section~6.3 the covariance is a $2\times2$ determinant of matrix elements of a transfer operator. Here the
% SOURCE: this sentence states the outcome of Theorem thm:rhcover, certified by the final ru2 run over [20, 300] (pk05/out/ru2_U2.out f6f2280c:
% VALID True, VERDICT COVERED for N = 3 and N = 4; see the SOURCE lines under the theorem). Route H was logged CLOSED (
% 2026-11-07T11:26:55Z).
% SOURCE (N = 4 on [300, 350]): the further ru2 run pk08/out/ru2_x2.out (b7452c92: VALID True l.61, VERDICT N=4 [300, 350]: COVERED
% l.63); see the SOURCE lines under the theorem.
operator is written in the hub frame of each column, whose axis is the middle spin, and the determinant is read as a matrix element of the second
exterior power of the operator. Pairs less than two transfer steps apart are evaluated directly. For all other pairs the leading modes of the exterior
power decide the sign, uniformly in the distance between the edges and in their distances to the ends.

\paragraph{The column transfer.}
% SOURCE: LEMMA_STRIP_RH_v1_3 section 0.
Let $\mathcal H$ be the space of O($N$)-invariant functions in $L^2(\omega^{\otimes3})$, put $u_1=\sigma_1\cdot\sigma_2$ and $u_3=\sigma_2\cdot\sigma_3$,
let $V$ be multiplication by $e^{J(u_1+u_3)}$, and let $K$ be the integral operator with kernel $\prod_{r=1}^3e^{J\sigma_r\cdot\sigma_r'}$. Then
$T_0=V^{1/2}KV^{1/2}$ is compact, self-adjoint and positive on $\mathcal H$, and $Z(P_3\times P_n)=\langle v,T_0^{n-1}v\rangle$ with $v=V^{1/2}\one$. It
has the same nonzero eigenvalues as the operator $T$ of Section~6.3. A horizontal edge of row $r$ between the columns $c$ and $c+1$ replaces the $c$-th
factor $K$ by its derivative $K_r$ in the coupling of row $r$. A vertical edge of column $c$ inserts multiplication by $s_e$ between the two factors
$V^{1/2}$ of column $c$, and two edges in one slot insert the product of their insertions.

\paragraph{The hub frame.}
% SOURCE: LEMMA_STRIP_RH_v1_3 R1 (Schur for O(N-1), N-1 in {2,3}: finding 2 of the review).
Fix $\sigma_2=e$. As in Section~6.4, the harmonics of degree $l$ in $\sigma_1$ split under the stabilizer O($N-1$) of $e$ into sectors $m=0,\dots,l$, and
those of sector $m$ are $f_{l,m}(u_1)Y_{m,j}(\omega_1)$, $1\le j\le D_m$, where $\omega_1$ is the direction of the part of $\sigma_1$ orthogonal to $e$
and the $Y_{m,j}$ are an orthonormal basis of the sector. For $N-1\in\{2,3\}$ every sector is absolutely irreducible, so by Schur's lemma an invariant
function of $(\sigma_1,\sigma_3)$ pairs equal sectors of $\sigma_1$ and $\sigma_3$ through the canonical pairing. Hence $\mathcal H$ has the orthonormal
basis
\[
\Phi_{l_1l_3m}=D_m^{-1/2}\sum_{j=1}^{D_m}f_{l_1,m}(u_1)Y_{m,j}(\omega_1)\,f_{l_3,m}(u_3)Y_{m,j}(\omega_3),
\]
with $0\le m\le\min(l_1,l_3)$, and multiplication by $u_1$ (by $u_3$) acts in sector $m$ as $\Jac^{(m)}$ on the label $l_1$ (on $l_3$).

\begin{lemma}\label{lem:rhframe}
Let $l_1,l_3\ge0$, $s=\min(l_1,l_3)+1$ and $\mathcal H_{l_1l_3}=\mathrm{span}\{\Phi_{l_1l_3m}:0\le m<s\}$.
\begin{enumerate}
\item[(i)] The Laplacian $\Delta_2$ in $\sigma_2$ preserves $\mathcal H_{l_1l_3}$, and $-\Delta_2$ acts there by the tridiagonal matrix $C=C^{(l_1,l_3)}$
with, for $N=3$,
\begin{gather*}
C_{mm}=l_1(l_1+1)+l_3(l_3+1)-2m^2,\\
C_{m,m+1}=-\sqrt{(l_1-m)(l_1+m+1)(l_3-m)(l_3+m+1)},
\end{gather*}
the off-diagonal entry multiplied by $\sqrt2$ when $m=0$, and, for $N=4$,
\begin{gather*}
C_{mm}=l_1(l_1+2)+l_3(l_3+2)-2m(m+1),\\
C_{m,m+1}=-2(m+1)\sqrt{\frac{((l_1+1)^2-(m+1)^2)\,((l_3+1)^2-(m+1)^2)}{(2m+1)(2m+3)}}.
\end{gather*}
\item[(ii)] The off-diagonal entries are nonzero, and the eigenvalues of $C$ are $l_2(l_2+N-2)$ for $l_2=|l_1-l_3|+2j$, $0\le j<s$, each simple.
\item[(iii)] Let $Q=Q^{(l_1,l_3)}$ be the orthogonal matrix of unit eigenvectors of $C$, with columns labelled by $l_2$ and $Q^{(l_3,l_1)}:=Q^{(l_1,l_3)}$.
The product singlets $\Psi_{l_1l_2l_3}=\sum_mQ_{m,l_2}\Phi_{l_1l_3m}$ form an orthonormal basis of $\mathcal H$, and
\[
K\Psi_{l_1l_2l_3}=k_{l_1}k_{l_2}k_{l_3}\Psi_{l_1l_2l_3},\qquad K_r\Psi_{l_1l_2l_3}=r_{l_r}k_{l_1}k_{l_2}k_{l_3}\Psi_{l_1l_2l_3},
\]
with $k_l=k_l(J)$ as in \eqref{eq:FH}, $r_l=\partial_J\log k_l(J)$ and $l_r$ the label of row $r$.
\end{enumerate}
\end{lemma}

\begin{proof}
% SOURCE: LEMMA_STRIP_RH_v1_3 R2, proof of the entries (the derivation, Steps 1-4; Step 3 with the note 1) and proof of the spectrum.
Take $e=e_N$ and write $L^{(i)}_{ab}$ for the generators of rotations of $\sigma_i$ and $B^{(i)}_a=L^{(i)}_{aN}$. On $\mathcal H$,
$L^{(1)}+L^{(2)}+L^{(3)}=0$, so $(L^{(2)}_{ab})^2=(L^{(1)}_{ab}+L^{(3)}_{ab})^2$ there, and
\[
-\Delta_2=-\Delta_1-\Delta_3-2\sum_{a<b<N}L^{(1)}_{ab}L^{(3)}_{ab}-2\sum_{a<N}B^{(1)}_aB^{(3)}_a .
\]
In sector $m$, $\sum_{a<b<N}L^{(1)}_{ab}L^{(3)}_{ab}=m(m+N-3)$, which gives the diagonal. The last sum maps sector $m$ to the sectors $m\pm1$. Write
$f_{l,m}(u)=c_{l,m}(1-u^2)^{m/2}p^{(m)}_{l-m}(u)$ with $c_{l,m}>0$ and $p^{(m)}_n$ orthonormal for the weight $(1-u^2)^{m+\lambda-1/2}$. The component in
sector $m+1$ comes from the three-term rule $tP_m=\alpha'_mP_{m+1}+\gamma'_mP_{m-1}$ of the zonal polynomials of $S^{N-2}$ and from the ladder
$(p^{(m)}_n)'=\kappa\,p^{(m+1)}_{n-1}$, $\kappa>0$, whose size the Sturm--Liouville identity of the Gegenbauer equation fixes. This gives
$C_{m,m+1}=-2\alpha'_m\sqrt{D_m/D_{m+1}}\,\rho_{l_1,m}\rho_{l_3,m}$ with $\rho_{l,m}^2=(l-m)(l+m+N-2)$, which are the stated entries. For $m<s-1$ they
are nonzero, so $C$ is an unreduced Jacobi matrix with $s$ simple eigenvalues. Each eigenvalue is $l_2(l_2+N-2)$ for a degree $l_2$ in $\sigma_2$ that
carries an invariant of $H_{l_1}\otimes H_{l_2}\otimes H_{l_3}$. Such an $l_2$ satisfies $l_2\le l_1+l_3$, because $\Phi_{l_1l_3m}$ is a polynomial of
degree $l_1+l_3$ in $\sigma_2$, and by symmetry in the three variables also $l_1\le l_2+l_3$ and $l_3\le l_1+l_2$; and $l_1+l_2+l_3$ is even, because
$-1\in\mathrm O(N)$ acts on $H_l$ by $(-1)^l$. There are exactly $s$ such $l_2$, so each occurs once. For (iii), \eqref{eq:FH} and \eqref{eq:proj} give
$K=k(\Delta_1)k(\Delta_2)k(\Delta_3)$ with $k(-l(l+N-2))=k_l$, and $\Delta_2$ commutes with $\Delta_1$, $\Delta_3$ and O($N$); $K_r$ replaces the
factor $k_{l_r}$ of row $r$ by $\partial_Jk_{l_r}$.
\end{proof}

\noindent The sign of the off-diagonal entries does not depend on choices of sign inside a sector. Computations may use the basis $(-1)^m\Phi_{l_1l_3m}$,
in which these entries are positive; this changes $Q$ by a diagonal of signs and none of the quantities below. For $s\ge2$ a column $Q_{\cdot,l_2}$, with
its sign fixed by $\langle Q_{\cdot,l_2},U\rangle>0$ and the same sign used for $(l_3,l_1)$, is enclosed from a computed vector $U$ by
$\|Q_{\cdot,l_2}-U\|\le\sqrt2\,\|CU-\mu U\|/(g\|U\|)+|\|U\|-1|$, with $\mu=l_2(l_2+N-2)$ and $g\ge4|l_1-l_3|+2N$ its distance to the other eigenvalues;
for $s=1$, $Q=(1)$.
% SOURCE: LEMMA_STRIP_RH_v1_3 R2 (gauge, Q radius, finding 7, s >= 2; for s = 1, Q = (1)).

\paragraph{Frames and truncation.}
% SOURCE: LEMMA_STRIP_RH_v1_3 R3 (normalization, (i)-(iv)) and R4 (insertions, (ii) frames, (iii) Davis-Kahan in its sector form, findings 9-11, N4).
Put $\hat k_l=c\,k_l$ with one constant $c>0$ for all $l$, $\tilde E_m=e^{J(\Jac^{(m)}-1)}$, $\hat E=\bigoplus_m\tilde E_m^{1/2}\otimes\tilde E_m^{1/2}$ (on
the labels $l_1$ and $l_3$), $\mathsf D=\mathrm{diag}(\hat k_{l_1}\hat k_{l_2}\hat k_{l_3})$ on the product singlets and $H=Q\mathsf DQ^\top$. Since
$\hat E=e^{-J}V^{1/2}$ and $H$ is a positive multiple of $K$, the operator $T=\hat EH\hat E$ is a positive multiple of $T_0$. With $S=H^{1/2}\hat E$ we
have $T=S^*S$, and we work mostly with $\tilde T=SS^*$, which has the same nonzero eigenvalues. For a cutoff $L$ let $H_L$ be the part of $H$ on the
blocks with $\max(l_1,l_3)\le L$, $S_L=H_L^{1/2}\hat E$, $T_L=S_L^*S_L$, and let
\[
\hat T=\mathsf D^{1/2}Q^\top GQ\,\mathsf D^{1/2},\qquad G=\bigoplus_m(\tilde E_m)_{[m,L]}\otimes(\tilde E_m)_{[m,L]},
\]
on the product singlets with $\max(l_1,l_3)\le L$, where $(\cdot)_{[m,L]}$ is the compression to the labels $m\le l\le L$. A horizontal edge $e$ of row
$r$ acts by $R_e=\mathrm{diag}(r_{l_r})$ on the product singlets. A vertical edge acts by $U_e$, multiplication by $u_1$ (tv) or by $u_3$ (bv); we put
$\tilde T_e=SU_eS^*$ and $\hat T_e=\mathsf D^{1/2}Q^\top G_eQ\,\mathsf D^{1/2}$, where $G_e=\bigoplus_m(\Jac^{(m)}\tilde E_m)_{[m,L]}\otimes(\tilde E_m)_{[m,L]}$
for tv and the two factors are exchanged for bv. The reflection $\rho$ that exchanges $\sigma_1$ and $\sigma_3$ acts by $\rho\Psi_{l_1l_2l_3}=\Psi_{l_3l_2l_1}$;
we write $\pm$ for its two eigenspaces, the sectors.

\begin{lemma}\label{lem:rhtrunc}
\begin{enumerate}
\item[(i)] $\|T-T_L\|\le\varepsilon_T:=\hat k_{L+1}\hat k_0^2$.
\item[(ii)] $T_L$ and $\hat T$ have the same nonzero eigenvalues, with multiplicities. If $\chi$, written in the product singlets, is a unit eigenvector
of $\hat T$ with eigenvalue $\nu>0$, then $S_L^*\chi/\sqrt\nu$ is a unit eigenvector of $T_L$. In the same coordinates $S_LU_eS_L^*$ is $\hat T_e$.
\item[(iii)] $T$, $T_L$ and $\hat T$ are positive and commute with $\rho$, and in each sector the $k$-th eigenvalues satisfy
$|\lambda_k^\pm(T)-\lambda_k^\pm(\hat T)|\le\varepsilon_T$, where $\lambda_k^\pm(\hat T)=0$ beyond the rank.
\item[(iv)] Let $g$ be the distance from an eigenvalue of $T_L$ to the other eigenvalues of $T_L$ in its sector, zero included. If $g>2\varepsilon_T$,
then $T$ has exactly one eigenvalue of that sector within $\varepsilon_T$ of it, and the unit eigenvectors, with signs aligned, differ by at most
$\sqrt2\,\varepsilon_T/(g-\varepsilon_T)$.
\end{enumerate}
\end{lemma}

\begin{proof}
The coefficients $\hat k_l$ decrease in $l$, since $k_{l+1}/k_l=I_{l+\lambda+1}(J)/I_{l+\lambda}(J)<1$ \cite[(16)]{Am74}. The part of $H$ on the blocks
with $\max(l_1,l_3)\ge L+1$ is positive, with norm at most $\hat k_{L+1}\hat k_0^2$, and $\|\hat E\|\le1$ because $\|\Jac^{(m)}\|\le1$; this gives (i). For
(ii), $S_LS_L^*=H_L^{1/2}P_L\tilde EP_LH_L^{1/2}$, where $\tilde E=\hat E^2$ and $P_L$ projects onto the labels $l_1,l_3\le L$; so $P_L\tilde EP_L$ is the direct
sum of the compressions, and conjugation by $Q$ gives $\hat T$. The same computation with $U_e\tilde E$ in place of $\tilde E$ gives $\hat T_e$: it is the
compression of the product $\Jac^{(m)}\tilde E_m$, not the product of compressions. For (iii), $\hat E$, $H$ and $H_L$ commute with $\rho$, because
$\mathsf D$ is symmetric in $l_1$ and $l_3$ and $Q^{(l_3,l_1)}=Q^{(l_1,l_3)}$, and Weyl's inequality applies in each sector. For (iv), Weyl's inequality
isolates the eigenvalue $\nu$ of $T$. If $\psi$ and $\psi^L$ are the eigenvectors and $P_\perp$ projects off $\psi^L$ inside the sector, then
$(T_L-\nu)P_\perp\psi=-P_\perp(T-T_L)\psi$, and every other eigenvalue of $T_L$ in the sector is at distance at least $g-\varepsilon_T$ from $\nu$, so
$\sin\angle(\psi,\psi^L)\le\varepsilon_T/(g-\varepsilon_T)$ and $\|\psi-\psi^L\|\le\sqrt2\sin\angle(\psi,\psi^L)$.
\end{proof}

\noindent The truncation thus enters once, through $\varepsilon_T$, in the frame of $T$, where it is linear, while $\hat T$ and $\hat T_e$ need only the
compressions of $\tilde E_m$ and $\Jac^{(m)}\tilde E_m$, which are enclosed with proven tails as in the proof of Theorem~\ref{thm:hubcover}.

\paragraph{Objects.}
% SOURCE: LEMMA_STRIP_R7R8_v1_1 R7 (R7a transfer formulas, R7b objects, same-gap pairs, the ends) and R8 (L2), centring.
Put $w=Sv$ and $w_p=\tilde T^pw$ for $p\ge0$. For a horizontal type $e$ put $w_{e,p}=R_ew_p$; for a vertical type $e$ put $w_{e,0}=SU_ev$ and
$w_{e,p}=\tilde T_ew_{p-1}$ for $p\ge1$. An edge has index $p$ from the left end if it lies between the columns $p+1$ and $p+2$ (horizontal) or in the
column $p+1$ (vertical); its index from the right end is defined in the mirror way, and the mirror symmetry of the strip gives the same vectors at the
right end. For $a,b,c,d\in\mathcal H$ put $\langle a\wedge b,c\wedge d\rangle=\langle a,c\rangle\langle b,d\rangle-\langle a,d\rangle\langle b,c\rangle$, and
let $\Lambda^2\tilde T$ act by $a\wedge b\mapsto\tilde Ta\wedge\tilde Tb$.

\begin{lemma}\label{lem:rhobj}
Let $e$ have index $p$ from the left end and $f$ index $q$ from the right end, with $k\ge0$ complete factors $\tilde T$ between their insertions, so that
$n=p+k+q+2$. Then $Z=\langle w_p,\tilde T^kw_q\rangle$ is a positive multiple of $Z(P_3\times P_n)$, and
\[
\Cov(s_e,s_f)\,Z^2=\big\langle w_p\wedge w_{e,p},(\Lambda^2\tilde T)^k(w_q\wedge w_{f,q})\big\rangle .
\]
The right side does not change if $w_{e,p}$ is replaced by $w'_{e,p}=w_{e,p}-\alpha_ew_p$, or $w_{f,q}$ by $w'_{f,q}=w_{f,q}-\alpha_fw_q$, for any real
$\alpha_e,\alpha_f$. Two horizontal edges of one junction are the case $k=0$. Every pair of distinct edges of $P_3\times P_n$, $n\ge4$, other than the two
vertical edges of one column, arises in this way.
\end{lemma}

\begin{proof}
For $m\ge1$, $T^m=S^*\tilde T^{m-1}S$ and $ST^m=\tilde T^mS$, so the partition function is a positive multiple of $\langle w,\tilde T^{n-2}w\rangle$. A
horizontal insertion turns one factor $T$ into $S^*R_eS$, because $R_e$ commutes with $H^{1/2}$; a vertical insertion in a column $c\ge2$ turns one
factor $\tilde T$ into $\tilde T_e$, and in the first column it gives $w_{e,0}$. Collecting the factors on both sides of the insertions gives $Z$,
$Z_e=\langle w_{e,p},\tilde T^kw_q\rangle$, $Z_f=\langle w_p,\tilde T^kw_{f,q}\rangle$ and $Z_{ef}=\langle w_{e,p},\tilde T^kw_{f,q}\rangle$ for the
derivatives of $Z$ in the couplings of $e$ and $f$, and $\Cov(s_e,s_f)Z^2=ZZ_{ef}-Z_eZ_f$ is the stated matrix element of $(\Lambda^2\tilde T)^k$. The
centring holds because $a\wedge(b-\alpha a)=a\wedge b$. For two horizontal edges of one junction the factor of that junction carries
$R_eR_f$, and $\langle w_p,R_eR_fw_q\rangle=\langle w_{e,p},w_{f,q}\rangle$. Two edges in different slots have their insertions in different factors;
the only two distinct edges in one slot other than a pair at one junction are the two vertical edges of one column.
\end{proof}

\paragraph{Spectral enclosures.}
% SOURCE: LEMMA_STRIP_RH_v1_3 R5 (i)-(v) (finding 13: x exactly in the sector) and R6 (leading mode, finding 16).
\begin{lemma}\label{lem:rhspec}
\begin{enumerate}
\item[(i)] For $x\ne0$ and real $\mu$, $\hat T$ has an eigenvalue within $\|\hat Tx-\mu x\|/\|x\|$ of $\mu$.
\item[(ii)] If disjoint intervals $I_1,\dots,I_K$ each contain an eigenvalue of $\hat T$ in the sector $\pm$, then every other eigenvalue of that sector is
at most $\tr(\hat T\Pi_\pm)-\sum_i\max(\min I_i,0)$, where $\Pi_\pm=(1\pm\rho)/2$, $\tr(\hat T\Pi_\pm)=(\tr\hat T\pm\tr\hat T\rho)/2$ and
\begin{gather*}
\tr\hat T=\sum\mathsf D_{l_1l_2l_3}\sum_mQ_{m,l_2}^2\,\tilde E_m(l_1,l_1)\tilde E_m(l_3,l_3),\\
\tr\hat T\rho=\sum\mathsf D_{l_1l_2l_3}\sum_mQ_{m,l_2}^2\,\tilde E_m(l_1,l_3)^2,
\end{gather*}
the outer sums running over the product singlets with $\max(l_1,l_3)\le L$.
\item[(iii)] Let $x$ lie in one sector, let $\nu$ be a simple eigenvalue of $\hat T$ in that sector with unit eigenvector $\chi_\nu$, and let every other
eigenvalue of the sector lie at distance at least $\gamma>0$ from $\mu$. Then $\sin\angle(x,\chi_\nu)\le\|\hat Tx-\mu x\|/(\gamma\|x\|)$.
\item[(iv)] The largest eigenvalue of $T$ is simple, its eigenvector $\psi_0$ is positive and $\rho$-even, and $\langle v,\psi_0\rangle>0$.
\end{enumerate}
\end{lemma}

\begin{proof}
For (i) and (iii), $\|(\hat T-\mu)x\|$ is at least the distance from $\mu$ to the spectrum times $\|x\|$, and at least $\gamma$ times the norm of the part of
$x$ orthogonal to $\chi_\nu$. (ii) holds because $\hat T\ge0$. The trace formulas hold because $G$ is diagonal in $m$ and $\rho$ maps the entry of $G$ at
$(l_1,l_3,m)$ to the entry at $(l_3,l_1,m)$, while $\tilde E_m$ is symmetric. For (iv), the kernel of $T_0$ on $(S^{N-1})^3$ is continuous and strictly
positive, so by Jentzsch's theorem its largest eigenvalue is simple with a positive eigenvector; by uniqueness this eigenvector is invariant under
O($N$) and under $\rho$, so it lies in $\mathcal H$; and $\langle v,\psi_0\rangle>0$ because $v>0$.
\end{proof}

\paragraph{Certified modes.}
% SOURCE: LEMMA_STRIP_RH_v1_3 R3 (iv), R4 (iii), R5, R6 (far-pair hypotheses, findings 14, 15, N1; (G)) and section 7; LEMMA_STRIP_R7R8_v1_1 R8
% (setting for the mode expansions; W6); LEMMA_STRIP_R9_v1_4 TR4 (mode vectors), TR6 (spectral data under truncation).
On a subinterval of $J$ the computation produces disjoint intervals that contain eigenvalues $\nu_0>\nu_1>\dots>\nu_5$ of $\tilde T$, each simple in its
sector, unit eigenvectors $\chi_0,\dots,\chi_5$ enclosed in angle, and a number $u_{\rm rem}$ that bounds every other eigenvalue of $\tilde T$ in both
sectors. They come from Lemma~\ref{lem:rhspec} for $\hat T$, with mode vectors that lie exactly in one sector, and are transferred to $\tilde T$ by
Lemma~\ref{lem:rhtrunc}: every interval and both sector bounds are widened by $\varepsilon_T$, and the angles carry the Davis--Kahan term of
Lemma~\ref{lem:rhtrunc}(iv), transported by $\chi_a=S\psi_a/\sqrt{\nu_a}$. The computation checks that the interval of $\nu_0$ lies above all other
bounds, that $\chi_1$ is $\rho$-even, and that
$u_{\rm rem}<\nu_1$.

\begin{lemma}\label{lem:rhinf}
Put $\theta_a=\langle\chi_a,w\rangle$ and $w_\perp=w-\sum_{a\le5}\theta_a\chi_a$, assume $\theta_0\ne0$, and for $P\ge1$ let
\[
\eta(P)=\sum_{a=1}^5\Big(\frac{\nu_a}{\nu_0}\Big)^P\Big|\frac{\theta_a}{\theta_0}\Big|+\Big(\frac{u_{\rm rem}}{\nu_0}\Big)^P\frac{|w_\perp|}{|\theta_0|}.
\]
\begin{enumerate}
\item[(i)] For every $p\ge P$, $w_p=\nu_0^p\theta_0(\chi_0+w^\circ_p)$ with $|w^\circ_p|\le\eta(P)$.
\item[(ii)] Let $e$ be horizontal, $R'=R_e-\alpha_e$ and $z_e=R'\chi_0-\langle\chi_0,R'\chi_0\rangle\chi_0$. The norms $\|R'\chi_a\|$ and $\|R'\tilde T\|$ are
finite, and for every $p\ge P$
\begin{gather*}
w_p\wedge w'_{e,p}=(\nu_0^p\theta_0)^2\big(\chi_0\wedge z_e+\Delta_e(p)\big),\\
|\Delta_e(p)|\le m_e+\eta(P)\|R'\chi_0\|+\eta(P)m_e,\\
m_e=\sum_{a=1}^5\Big(\frac{\nu_a}{\nu_0}\Big)^P\Big|\frac{\theta_a}{\theta_0}\Big|\,\|R'\chi_a\|+\frac{\|R'\tilde T\|\,u_{\rm rem}^{P-1}|w_\perp|}{\nu_0^P|\theta_0|}.
\end{gather*}
\item[(iii)] Let $e$ be vertical and $T'=\tilde T_e-\alpha_e\tilde T$, so that $w'_{e,p}=T'w_{p-1}$. For every $p\ge P\ge2$,
$w_p\wedge w'_{e,p}=\nu_0(\nu_0^{p-1}\theta_0)^2(\chi_0\wedge z_e+\Delta_e(p))$, where $z_e$ is the part of $T'\chi_0$ orthogonal to $\chi_0$ and
$\Delta_e(p)$ obeys the bound of (ii) with $T'$ in place of $R'$ and $P-1$ in place of $P$ in $m_e$.
\end{enumerate}
\end{lemma}

\begin{proof}
% SOURCE: LEMMA_STRIP_R7R8_v1_1 R8 INF1-INF3 (W1: ||R'T~|| in place of the infinite ||R'||) and LEMMA_STRIP_R9_v1_4 TR7.
(i) $w_p=\sum_{a\le5}\nu_a^p\theta_a\chi_a+\tilde T^pw_\perp$, and $\tilde T$ maps the orthogonal complement of $\chi_0,\dots,\chi_5$ into itself with norm at
most $u_{\rm rem}$; every term of $w^\circ_p$ decreases in $p$. (ii) The operator $R'$ is unbounded, since $r_l$ grows like $l/J$, but
$\chi_a=\tilde T\chi_a/\nu_a$ lies in the range of $H^{1/2}$, on which $R'$ is bounded because $\hat k_l$ decays faster than any geometric sequence, and
$\|R'\tilde T\|\le\|R'S\|\,\|S^*\|$. Since $R'w^\circ_p=\sum_{a\ge1}(\nu_a/\nu_0)^p(\theta_a/\theta_0)R'\chi_a+(R'\tilde T)\tilde T^{p-1}w_\perp/(\nu_0^p\theta_0)$,
we have $\|R'w^\circ_p\|\le m_e$ for $p\ge P$. Expanding $w_p\wedge R'w_p$ into four terms and using $|a\wedge b|\le|a||b|$ gives the bound; the projection
in $z_e$ does not change the wedge with $\chi_0$. (iii) is the same computation with $w_p=\nu_0\cdot\nu_0^{p-1}\theta_0(\chi_0+w^\circ_p)$ and
$T'w_{p-1}=\nu_0^{p-1}\theta_0(T'\chi_0+T'w^\circ_{p-1})$, where $T'$ is bounded.
\end{proof}

\noindent Let $D_e$ be a bound for $|\Delta_e(p)|$, $p\ge P_e$, as in Lemma~\ref{lem:rhinf}, where $P_e=9$ for mh and $P_e=6$ for the other types. Since
$\|\Lambda^2\tilde T\|=\nu_0\nu_1$, for every $\Upsilon=x\wedge y$ and every $p\ge P_e$ the value $\langle w_p\wedge w'_{e,p},(\Lambda^2\tilde T)^k\Upsilon\rangle$
is a positive multiple of $\langle\chi_0\wedge z_e,(\Lambda^2\tilde T)^k\Upsilon\rangle+E$ with $|E|\le D_e(\nu_0\nu_1)^k|x||y|$. One check of the representative
$\chi_0\wedge z_e$ with this error therefore covers every index $p\ge P_e$; we call $(e,\infty)$ the infinity object of type $e$. When both objects of a
pair are infinity objects, the product of their two errors is added. In the computation $\chi_0$ is replaced by its computed unit direction, and $D_e$
includes the resulting error.
% SOURCE: LEMMA_STRIP_R7R8_v1_1 R8 INF4 (W4: exact positive factor on the left, upper bounds only in the error) and the computed representatives
% (W1-W3); LEMMA_STRIP_R9_v1_4 TR7 (the representative's error sqrt2 s_0' (ne_0 + Rop_L); no term of order sqrt(eps_T)).

\paragraph{Near and far pairs.}
% SOURCE: LEMMA_STRIP_R7R8_v1_1 R8 (NEAR, FAR, monotone bound; W5, W6) and LEMMA_STRIP_RH_v1_3 section 7 (monotonicity).
For object pairs $\Xi=x\wedge y$ and $\Upsilon=x'\wedge y'$, with $x,x'$ plain and $y,y'$ centred inserted objects, put
$S(k)=\langle\Xi,(\Lambda^2\tilde T)^k\Upsilon\rangle$. For $k<2$ it is computed directly as a $2\times2$ determinant of inner products of the objects and of
their images under $\tilde T$. For $k\ge2$ the leading modes of $\Lambda^2\tilde T$ decide its sign.

\begin{lemma}\label{lem:rhnf}
Put $X_{ab}=\langle\chi_a\wedge\chi_b,\Xi\rangle$ and $Y_{ab}=\langle\chi_a\wedge\chi_b,\Upsilon\rangle$ for $0\le a<b\le5$, and let $\Xi^\perp$, $\Upsilon^\perp$ be
the parts of $\Xi$, $\Upsilon$ orthogonal to all $\chi_a\wedge\chi_b$ with $a<b\le5$. Then for $k\ge0$
\[
\frac{S(k)}{(\nu_0\nu_1)^k}=X_{01}Y_{01}+\sum_{(a,b)\ne(0,1)}r_{ab}^k\,X_{ab}Y_{ab}+R_k,\qquad|R_k|\le q^k|\Xi^\perp||\Upsilon^\perp|,
\]
where $r_{ab}=\nu_a\nu_b/(\nu_0\nu_1)<1$ for $(a,b)\ne(0,1)$ and $q=u_{\rm rem}/\nu_1<1$. If $\ell$ is a lower bound for $X_{01}Y_{01}$, then
$g(k)=\ell-\sum_{(a,b)\ne(0,1)}r_{ab}^k|X_{ab}Y_{ab}|-q^k|\Xi^\perp||\Upsilon^\perp|$ is nondecreasing in $k$ and $g(k)\le S(k)/(\nu_0\nu_1)^k$; so $g(k_*)>0$
gives $S(k)>0$ for every $k\ge k_*$. Moreover $|\Xi^\perp|\le|x_C||y_\perp|+|x_\perp||y|$, where $x_C$ and $x_\perp$ are the parts of $x$ in and orthogonal
to $\mathrm{span}\{\chi_0,\dots,\chi_5\}$, and likewise for $y$.
\end{lemma}

\begin{proof}
$\Lambda^2\tilde T$ is compact, self-adjoint and positive, with eigenvectors $\chi_a\wedge\chi_b$, $a<b$, and eigenvalues $\nu_a\nu_b$, where the $\chi_a$ run
over an eigenbasis of $\tilde T$ that extends $\chi_0,\dots,\chi_5$. A pair with an index beyond $5$ has $\nu_a\nu_b\le\nu_0u_{\rm rem}$, so the
Cauchy--Schwarz inequality bounds the sum over these pairs by $(\nu_0u_{\rm rem})^k|\Xi^\perp||\Upsilon^\perp|$. Each $r_{ab}$ and $q$ is less than one, so
every subtracted term in $g$ decreases in $k$. Finally $x\wedge y-x_C\wedge y_C=x_C\wedge y_\perp+x_\perp\wedge y$.
\end{proof}

\noindent For an infinity object the errors of the preceding paragraph, divided by $(\nu_0\nu_1)^k$, are bounded uniformly in $k$ and are subtracted in
$g$ as well. If $\chi_1$ were $\rho$-odd, then $X_{01}$ would vanish for every $\rho$-even object pair, and the check could only fail.

\begin{lemma}\label{lem:rhnp}
\begin{enumerate}
\item[(NORM)] If every object is replaced by a positive multiple, where the factor may depend on $J$ and on the number of factors of the object, then
every value $S(k)$ and every amplitude $X_{ab}$ is multiplied by a positive number. The same holds if $\tilde T$ is replaced by a positive multiple,
which does not change the ratios $\nu_a/\nu_b$.
\item[(PAR)] $\rho$ commutes with $\tilde T$, $\hat T$, $R_{\rm mh}$ and with scalars, exchanges $R_{\rm th}$ with $R_{\rm bh}$ and $\tilde T_{\rm tv}$ with
$\tilde T_{\rm bv}$, and fixes $w$. Hence the plain objects and all objects of type mh are $\rho$-even, and $\langle\chi_a,x\rangle=0$ whenever $\chi_a$ is
$\rho$-odd and $x$ is $\rho$-even.
\end{enumerate}
\end{lemma}

\begin{proof}
% SOURCE: LEMMA_STRIP_RH_v1_3 section F (NORM; PAR (a)-(c), V6).
(NORM) Each product of two inner products in $S(k)$ contains one plain and one inserted object from each side, and $X_{ab}$ is bilinear in the plain and
the inserted object. (PAR) $\mathsf D$ and $G$ are symmetric under $l_1\leftrightarrow l_3$, the multiplier $r_{l_2}$ of $R_{\rm mh}$ depends only on the middle
label, and $w=e^JH^{1/2}\tilde E\one$, where $\tilde E\one$ lies in sector $m=0$ and is symmetric in $l_1$ and $l_3$; $H^{1/2}=Q\mathsf D^{1/2}Q^\top$ commutes
with $\rho$ by the convention $Q^{(l_3,l_1)}=Q^{(l_1,l_3)}$. For $\rho$-odd $\chi_a$ and $\rho$-even $x$, $\langle\chi_a,x\rangle=\langle\rho\chi_a,\rho x\rangle=
-\langle\chi_a,x\rangle$.
\end{proof}

\paragraph{The computed frame.}
% SOURCE: LEMMA_STRIP_R9_v1_4 sections 1-3 (frames, insertion factors f_x (Y1), top interval (Y3), TR1 with M_e = max(1, alpha_e), TR2-TR5) and
% LEMMA_STRIP_RH_v1_3 R3 (v), R4 (iii) (eps_T, eps_e on a subinterval).
The computation works on closed subintervals $b_-\le J\le b_+$, $b_\pm=b_0\pm h$, with $t=J-b_0$ and the normalization $\hat k_l=e^{-t}k_l(J)/k_0(b_0)$.
Here $\hat k_0$ decreases in $t$, and $I_{l+\lambda}(J)/I_\lambda(J)$, a product of ratios $I_{\nu+1}/I_\nu$, increases in $J$ (proof of
Lemma~\ref{lem:rmono}). So, with $A=\sup\hat k_0$ and $\hat k_{0,\rm lo}=\inf\hat k_0$ over the subinterval,
\begin{gather*}
\varepsilon_T\le A^3\frac{I_{L+1+\lambda}(b_+)}{I_\lambda(b_+)},\\
\varepsilon_e:=A^3\frac{I_{L+\lambda}(b_+)}{I_\lambda(b_+)},\\
\varepsilon_W:=\max\Big(\frac{2A\varepsilon_T}{\hat k_{0,\rm lo}},\Big(1+\frac{L+1}{b_-}\Big)\varepsilon_e\Big),
\end{gather*}
and we put $\varepsilon_*=\max(\varepsilon_T,\varepsilon_e,\varepsilon_W)$. Every object is formed in the truncated frame, with $\hat T$, $\hat T_e$ and $R_e$ on
the labels $\max(l_1,l_3)\le L$ in place of $\tilde T$, $\tilde T_e$ and $R_e$; for an object $x$ of the exact frame we write $x_L$ for the computed one and
$p_x$ for its number of factors $T$ or $S^*R'S$. Let $[\nu_0^-,\nu_0^+]$ be the interval of $\nu_0$, widened by $\varepsilon_T$, and $|\theta_0|^-$ a lower
bound for $|\langle\chi_0^L,w_L\rangle|$, with $\chi_0^L$ the top eigenvector of $\hat T$. The insertion factor $f_x$ of an object is $1$ for plain objects,
$\max(\max_l|r_l-\alpha_e|,1)$, the maximum over the
computed labels, for horizontal ones, and $1+|\alpha_e|$ for vertical ones.

\begin{lemma}\label{lem:rhslack}
Assume $0\le\alpha_e\le1$ for the horizontal types and $p_{\max}^2\varepsilon_*/\nu_0^-\le10^{-6}$, where $p_{\max}=\max_xp_x$. Put
$c=\max_x\big((\nu_0^+/\nu_0^-)(1+\varepsilon_T/\nu_0^-)\big)^{p_x}$, $\kappa=|v|\sqrt{\nu_0^+}/|\theta_0|^-$ and $\varrho=\max_xf_x|w_{p_x,L}|/|x_L|^-$, where
$|x_L|^-$ is a lower bound for $|x_L|$; a mode vector $\chi_a$ enters with $p_x=0$, $x_L:=P_L\chi_a$ and $(\nu_0^+/\nu_a^-)^{1/2}$ in place of
$f_x|w_{p_x,L}|$. Then for all objects $x,y$, the
mode vectors and the representatives of the infinity objects included,
\[
|\langle x,y\rangle-\langle x_L,y_L\rangle|\le\varepsilon_{\rm sl}|x_L||y_L|,\qquad
\varepsilon_{\rm sl}=\frac{\varepsilon_*}{\nu_0^-}\big(2p_{\max}c\kappa\varrho+c^2\kappa^2\varrho^2\big)(1+10^{-6}).
\]
\end{lemma}

\begin{proof}
Each difference between an exact and a truncated operator, $T-T_L$, $S^*R'S-S_L^*R'S_L$ and $S^*R'R''S-S_L^*R'R''S_L$, is $\hat E$ times the tail of an
operator that is diagonal on the product singlets times $\hat E$. Its norm is at most the largest tail entry, a product of three factors $\hat k_{l_r}$ in
which an inserted row carries $r_{l_r}-\alpha_e$ once or twice. We use $|r_l-\alpha_e|\le\max(r_l,1)$, $r_l\hat k_l\le\hat k_{l-1}$ and $r_l<1+l/J$, and that
both $\hat k_l$ and $t_l=(1+l/J)\hat k_{l-1}$ decrease in $l$: for $t_l$, $l\ge1$, the Bessel ratio bound $I_{\nu+1}(J)/I_\nu(J)<J/(\nu+\frac12+J)$ from the
proof of Lemma~\ref{lem:rmono} and $\lambda\ge\frac12$ give $t_{l+1}/t_l<(J+l+1)J/(J+l)^2<1$. So these entries are at most $\varepsilon_T$,
$\varepsilon_e$ and $\varepsilon_W$, and the tails satisfy
$\|P_\perp S\|^2\le\varepsilon_T$ and $\|R'P_\perp S\|^2\le\varepsilon_W$, where $P_\perp$ projects onto the labels with $\max(l_1,l_3)>L$. Telescoping the
word of an object, with the computed factors to the left of each difference, gives $|P_Lx-x_L|\le p_x(\varepsilon_*/\nu_0^-)\kappa c_xf_x|w_{p_x,L}|$
with $c_x=((\nu_0^+/\nu_0^-)(1+\varepsilon_T/\nu_0^-))^{p_x}$, and the tail gives $|P_\perp x|\le(\varepsilon_*/\nu_0^-)^{1/2}\kappa c_xf_x|w_{p_x,L}|$; the factor
$\kappa$ converts $(\nu_0^+)^{p_x}|v|$ into the norm of the computed plain object, since $|w_{p,L}|\ge(\nu_0^-)^p|\theta_0|^-$. Then
$\langle x,y\rangle=\langle P_Lx,P_Ly\rangle+\langle P_\perp x,P_\perp y\rangle$ gives the claim. For a mode vector $x_L=P_L\chi_a$, so only its tail enters:
$\chi_a=S\psi_a/\sqrt{\nu_a}$ and $\|P_\perp S\|^2\le\varepsilon_T$ give $|P_\perp\chi_a|\le(\varepsilon_T/\nu_a^-)^{1/2}\le(\varepsilon_*/\nu_0^-)^{1/2}\kappa\,
(\nu_0^+/\nu_a^-)^{1/2}$, since $\kappa\ge1$: indeed $|\theta_0|^-\le|w_L|\le\|S_L\|\,|v|\le\sqrt{\nu_0^+}\,|v|$.
% SOURCE: the review AUDIT_RH_MODULE_v1_rv01_ADD1 (fc400fc4), M1 row: the one-clause reason for kappa >= 1, added at the fold-in.
\end{proof}

\noindent The coefficient of an object $x$ along a certified mode is enclosed from a computed mode vector $u_a$ by
\[
\langle\chi_a,x\rangle\in\frac{\langle u_a,x_L\rangle}{|u_a|}\pm\sqrt2\,s_a|x_L|\pm\varepsilon_{\rm sl}|x_L|,
\]
where $\sqrt2\,s_a$ bounds $|P_L\chi_a-u_a/|u_a||$. The number $s_a$ is the Krylov--Weinstein and Davis--Kahan bound for the angle between $u_a$ and the
eigenvector of $\hat T$, plus the transport of Lemma~\ref{lem:rhtrunc}(iv) to $\tilde T$, namely
$(\nu_0^+/\nu_a^-)^{1/2}\big(\sqrt2\,\varepsilon_T/(g_a-\varepsilon_T)+\varepsilon_T/(2\nu_a^-)\big)/\sqrt2$ with $g_a$ the gap of Lemma~\ref{lem:rhtrunc}(iv)
for $\nu_a$.
% SOURCE: LEMMA_STRIP_R9_v1_4 TR4 (the transport d_a) and TR6 (the enclosure of the mode coefficients); the review audit M1 and c3.

\noindent Vectors depending on $J$ are carried as Taylor models in $t$ of degree four, with exact floating-point coefficients and a rigorous bound for the
remainder in the $2$-norm; each operator step, multiplication by a polynomial in $t$ and linear combination adds its truncation and rounding errors to
the remainder. The inner products of the coefficient vectors come from one Gram matrix product in floating point, formed by a classical algorithm (not a
Strassen-type one), whose error is at most
$\gamma_{n+2}\|x\|\|y\|+(n+2)\eta$ for $n$ states and the underflow unit $\eta$, and the resulting polynomials in $t$ are evaluated on subballs of the
subinterval in ball arithmetic. With $\varepsilon_{\rm sl}|x_L||y_L|$ added to every inner product, every computed value of Lemmas~\ref{lem:rhinf}
and~\ref{lem:rhnf}, and every direct value $S(k)$ for $k<2$, has a lower end that bounds the true value from below on the subball.
% SOURCE: LEMMA_STRIP_R9_v1_4 section 5 (TM1-TM4; TM3 assumption Y7: classical gemm/syrk, single-threaded OpenBLAS) and TR9.

\begin{theorem}[$\dagger$]\label{thm:rhcover}
For $N=3$ and every $J\in[20,300]$, for $N=4$ and every $J\in[20,350]$, every $n\ge4$ and every pair of distinct edges $e,f$ of $P_3\times P_n$
other than the two vertical edges of one column, $\Cov(s_e,s_f)>0$.
\end{theorem}
% SOURCE (certificate): the final ru2 run over [20, 300] on the concatenated job lists of k1-k5, pk05/out/ru2_U2.out (f6f2280c;
% pk05/MANIFEST_ru2_U2_out.sha256 0979d81b, 4 OK): JOBS 358, counted 358; CONTROLS P-HOLE, P-VOIDJOB, P-ALPHA, P-TAG FIRED for N = 3 and N = 4;
% UNION one certified component and 0 gaps per N; VALID True; VERDICT N=3 [20, 300]: COVERED and N=4 [20, 300]: COVERED (l.361-367). The
% read: 910 intervals, 0 uncertified, 0 CONFLICT, engine b0f40c89 on every job. Batch k5's ru2: ru2_k5.out
% 820e6b1b (MANIFEST_ru2_k5_out fa7d57a5). Log, a live file and not the record: pk05/RU2_LOG_pk05.md l.48-52 (the entry
% "FINAL U2 [20, 300]" of 2026-11-07T11:23:24Z); the record is ru2_U2.out with its manifest.
% SOURCE: engine cs33h_1_rh_cert.py (b0f40c89) with core cs33b_rh_core.py (929ab8a0) and hubcert_v1e.py (c08d33c1); code map LEMMA_STRIP_R10_v1_4
% (bd7f8861); ledger and union LEMMA_STRIP_R11_v1_4 (5e532179); union check ru2_ledger_union.py (837a50b3).
% SOURCE (final numbers, not stated in the text): counted subintervals per part, from the JOB lines of ru2_U2.out, the same for N = 3 and N = 4:
% A 218 (k1-k5), B 199 (k1-k4), Bh 19 and Bv 19 (k5); 910 in all over both N. Smallest NEAR margin 6.234e-3 (N = 4, b0 59.4906, k1); smallest
% FAR margin 1.758e-7 (N = 3, top subinterval b0 300.8088, lead 7.523e-7, err 6.144e-7) (the read). Compute: 38.97 core-h of the 46 core-h cap.
% SOURCE (N = 4 on [300, 350]): pk08/R11_ADDENDUM_N4_300_350_pk08_v1_1.md (3797ca7e): one further ru2 run, pk08/out/ru2_x2.out (b7452c92;
% pk08/MANIFEST_ru2_x2_out.sha256 ff8f825c), on the job manifest pk08/MANIFEST_ru2_x2_jobs.sha256 (53f0db26; the chain's intervals on
% [302.31, 350] and batch k5's top N = 4 outputs, which overlap them): JOBS l.60, CONTROLS fired at both N l.59, VALID True
% l.61, VERDICT N=4 [300, 350]: COVERED l.63; the same engine cs33h_1_rh_cert.py (b0f40c89). With U2's VERDICT N=4 [20, 300] the
% two join at 300. N = 3 stays [20, 300].

\begin{proof}[Method]
% SOURCE: LEMMA_STRIP_R7R8_v1_1 R8 (coverage; W5; R9 TR8 on infinity pairs), LEMMA_STRIP_R10_v1_4 (controls and plants), LEMMA_STRIP_R11_v1_4
% (parts, union, counting rules, single engine U1, superseded batches).
By Lemma~\ref{lem:rhobj} and (NORM) it suffices to prove $S(k)>0$ for every pair of objects, where an object is $(e,p)$ with $p<P_e$ or the infinity object
$(e,\infty)$, which stands for every $p\ge P_e$. The interval $[20,300]$ ($[20,350]$ for $N=4$) is covered by closed subintervals. On each of them, separately for $N=3$ and
$N=4$, the following is certified.
\begin{enumerate}
\item[(a)] Inputs: Taylor models of degree four in $J$ with rigorous remainders for $\hat k_l$, $\hat k_l^{1/2}$, $r_l$, $\tilde E_m$ and $\Jac^{(m)}\tilde E_m$,
built with the enclosures used for Theorem~\ref{thm:hubcover}, and the blocks $Q$ with the radius bound above.
\item[(b)] Spectral data: the modes $\chi_0,\dots,\chi_5$ and $u_{\rm rem}$, with the checks listed above and the gap condition of
Lemma~\ref{lem:rhtrunc}(iv); the coefficients that vanish by (PAR) are set to zero.
\item[(c)] Objects: $w_p$ and $w'_{e,p}$ for $p<P_e$, the representatives of the infinity objects with their errors $D_e$, and the images of all of these
under $\tilde T$. Every propagation step is multiplied by a polynomial $\varpi(t)$ close to $1/\nu_0$, with a rigorous lower bound $\varpi>0$ on the
subinterval, which (NORM) allows; upper bounds of $\varpi$ enter only the error terms.
\item[(d)] Checks: for every unordered pair of objects and every $k\ge0$ with $p+q+k\ge2$, the direct value $S(k)$ for $k<2$. For $k\ge2$, the lower
bound for $S(k)/(\nu_0\nu_1)^k$ from the expansion of Lemma~\ref{lem:rhnf} must be positive at $k=2,\dots,k_*-1$, and the monotone bound must satisfy
$g(k_*)>0$ for some $k_*\le60$, which covers every $k\ge k_*$. An infinity object paired with itself stands for all indices
$p,q\ge P_e$, unequal ones included, and is checked like every other pair. The pairs at $k=0$ of one horizontal type, which are not pairs of distinct
edges, are checked as well.
\item[(e)] Controls: on the first subinterval of each run, deliberately falsified inputs must each make a check fail. Every part has the common ones, a
test of the ball arithmetic and a perturbed Casimir block, operator enclosure, trace bound, cutoff and infinity error; part A adds a perturbed mode
vector and a negated leading amplitude, and part B (or Bh and Bv) a negated near value. In every run the floating-point spectral data must agree with
an independent computation.
\end{enumerate}
The checks are divided into part A, the checks for $k\ge2$, and part B, the checks for $k<2$; part B may also be run as two parts Bh and Bv, which
together contain its checks. Each part carries the common controls and those of its own checks. A subinterval counts for a part when every control and
every check of the part holds on it.
With $U_P$ the union of the closed subintervals that count for part $P$, a range of $J$ is covered when it lies in $U_A\cap(U_B\cup(U_{Bh}\cap U_{Bv}))$.
A run counts only if it was made with the single production code, which its output identifies by checksum, if every control fired, if it prints its
centring constants at least as often as it certifies subintervals, and if every printed centring constant $\alpha_e$ of a horizontal type lies in
$[0,1-10^{-9}]$, so that Lemma~\ref{lem:rhslack} applies. An independent program, which
shares no code with the certification, forms the unions in exact rational arithmetic from the hexadecimal endpoints of the subintervals and reports
every gap; it has its own controls (a removed subinterval, a run without valid controls, a relabelled part and a centring constant out of range must
each be detected). Its final run over $[20,300]$, on the runs of all batches together, reports no gap for $N=3$ and for $N=4$, and a further run over $[300,350]$ for $N=4$ reports none either.
% SOURCE: the final run, ru2_U2.out (f6f2280c), reports one certified component and 0 gaps for N = 3 and N = 4 (UNION lines, l.361-362), all
% eight controls fired (l.363) and VALID True (l.365).
% SOURCE (N = 4 on [300, 350]): ru2_x2.out (b7452c92) reports one certified component and 0 gaps for N = 4 (UNION N=4, l.58), all
% eight controls fired (l.59) and VALID True (l.61).
% SOURCE (batches, R11): [20, 60] (q = 2%), [60, 150], [150, 200], [200, 250], [250, 300] (q = 1%), tags k1-k5 on cs33h_1 only, each closed by ru2
% on its job list before the next; superseded and ignored by ru2: cs33e certificates, the cs33b-cs33h test outputs, batches g1 (cs33g) and h1 (cs33h).
% SOURCE (counts and margins): see the SOURCE lines under Theorem thm:rhcover (910 counted subintervals; NEAR 6.234e-3, FAR 1.758e-7; 38.97 core-h).
\end{proof}

\begin{corollary}\label{cor:rhstrips}
For $N=3$ and every $J\in[20,300]$, for $N=4$ and every $J\in[20,350]$, and every $n\ge1$, $\Cov(s_e,s_f)\ge0$ for every pair of distinct edges
of $P_3\times P_n$. For $n\ge4$ the inequality is strict for every pair other than the two vertical edges of one column.
\end{corollary}

\begin{proof}
% SOURCE: LEMMA_STRIP_RH_v1_3 section 0 (scope, finding 1 and N3; V1, V2) and section 7 (engine scope).
For $n\ge4$ this is Theorem~\ref{thm:rhcover}, together with Theorem~\ref{thm:shared} for the two vertical edges of one column, which share the middle
vertex. For $n=3$ it is Theorem~\ref{thm:D}(i). For $n=2$ the graph $P_3\times P_2$ is the ladder with three rungs, a $2$-connected outerplanar graph
whose weak dual is a path: Theorem~\ref{thm:B}(a) applies at $N=4$, and Theorem~\ref{thm:C} at $N=3$, since an outerplanar graph has no $K_4$ minor,
hence no $W_4$ minor, and $G\cup\{e,f\}=G$ for edges $e,f$ of $G$. For $n=1$ the graph is a path, whose two edges lie in different blocks, so
$\Cov(s_e,s_f)=0$ by Lemma~\ref{lem:blocks}.
\end{proof}

\noindent Theorem~\ref{thm:D}(ii) covers $0<J\le20$ for $n\ge3$, and the arguments for $n\le3$ in the proof above hold at every $J\ge0$. With
Corollary~\ref{cor:rhstrips}, therefore,
$\Cov(s_e,s_f)\ge0$ holds for every pair of distinct edges of every $P_3\times P_n$ at every $0<J\le300$ for $N=3$ and every $0<J\le350$ for $N=4$.
% SOURCE: LEMMA_STRIP_R11_v1_4 (U4): certint6 tops 20.06497 (N = 4) and 20.07052 (N = 3), UNIFORM_0_20_AUDIT_v1 (e7f6c283).
